\documentclass{../notatki}

\title{Gry Kombinatoryczne}

\begin{document}

\section{Gra}

W ramach przedmiotu będziemy rozważać gry skończone, z pełną informacją o stanie
gry. Zakładamy, że w ramach gier będą występować tylko dwaj gracze, i co za tym
idzie, gra się może kończyć z wynikiem 1 (gracz 1 wygrywa), 0 (gracz 2 wygrywa),
lub 1/2 (remis).

\subsection{Drzewo gry}

Różne stany gry można rozważać w postaci drzewa gry, gdzie każdy węzeł
reprezentuje możliwy stan gry, krawędzie ruchy, a poziomy drzewa reprezentują
rundy gry.

\subsection{Strategia}

Strategią, nazywamy funckję działającą na podstawie stanu gry i zwracającą
ruch gracza, lub kolejny stan.

\textbf{W każdej grze ktoś ma strategię nieprzegrywającą}.

\textbf{Jeśli nie ma remisu w grze, to ktoś ma strategię wygrywającą}

\section{Analiza wstecz}

Analizując wstecz zachowania graczy w ramach gry, możemy przyporządkować
każdemu stanowi "wartość", która odpowiada za wynik gry w tym
momencie, zakładając optymalne ruchy obu graczy.

\begin{figure}[!htb]
  \centering
  \begin{tikzpicture}
    \node[circle, draw, fill=blue, minimum size=1cm] (parent) at (0,0) {1/2};

    \node[circle, draw, fill=red, minimum size=1cm, below left=of
    parent] (child1) {0};
    \node[circle, draw, fill=red, minimum size=1cm, below
    right=of parent] (child2) {1/2};

    \node[circle, draw, fill=blue, minimum size=1cm, below left
    = of child1] (grandchild1) {1};
    \node[circle, draw, fill=blue, minimum size=1cm, below
    = of child1] (grandchild2) {0};
    \node[circle, draw, fill=blue, minimum size=1cm, below
    = of child2] (grandchild3) {1};
    \node[circle, draw, fill=blue, minimum size=1cm, below right
    = of child2] (grandchild4) {1/2};

    \node[below left= of grandchild1] (grandgrandchild1) {1/2};
    \node[below= of grandchild1] (grandgrandchild2) {1};
    \node[below= of grandchild2] (grandgrandchild3) {0};
    \node[below left= of grandchild3] (grandgrandchild4) {1};
    \node[below= of grandchild3] (grandgrandchild5) {0};
    \node[below right= of grandchild4] (grandgrandchild6) {1/2};
    \node[below right= of grandchild4] (grandgrandchild7) {1/2};
    \node[below= of grandchild4] (grandgrandchild8) {0};

    \draw (parent) -- (child1);
    \draw (parent) -- (child2);
    \draw (child1) -- (grandchild1);
    \draw (child1) -- (grandchild2);
    \draw (child2) -- (grandchild3);
    \draw (child2) -- (grandchild4);

    \draw (grandchild1) -- (grandgrandchild1);
    \draw (grandchild1) -- (grandgrandchild2);
    \draw (grandchild2) -- (grandgrandchild3);
    \draw (grandchild3) -- (grandgrandchild4);
    \draw (grandchild3) -- (grandgrandchild5);
    \draw (grandchild4) -- (grandgrandchild6);
    \draw (grandchild4) -- (grandgrandchild7);
    \draw (grandchild4) -- (grandgrandchild8);
  \end{tikzpicture}
  \caption{Przykładowe wartościowanie stanów gry \label{fig:gry}}
\end{figure}

W kontekście \cref{fig:gry}, na podstawie analizy wstecz możemy bardzo łatwo
określić, że wartością gry jest 1/2. Na podstawie tego wynika, że każdy z graczy
ma strategię nieprzegrywającą. Równocześnie nikt nie ma strategii wygrywającej.
W kontekście takiego drzewa, strategią jest decyzja na każdym liściu, do którego
dziecka przejść.

\section{Hex}

Jeśli na każdym polu planszy Hex $n \times n$ położymy pionka o
dowolnym kolorze, to jeden z kolorów będzie w pozycji wygrywającej. Stąd też,
nie ma remisów w tej grze.

\section{Kradzież strategii}

Jest to technika analizy gry, która polega na analizie potencjalnych strategii
przeciwnika. Na przykład, pozwala nam udowodnić, że gracz ma strategię
wygrywającą, nawet jeśli nie jest w stanie ją znaleźć samodzielnie.

W Hex, nie ma remisów. W związku z tym, jak zakładamy, że gracz II ma strategię
wygrywyającą, to gracz I musi nie mieć wygrywającej strategii. Można, jednak
znaleźć sprzeczność w założeniach. Wyobraźmy sobie, że gracz I ma strategię
gracza II. W takiej sytuacji, gracz I może zrobić jakiś losowy ruch, i potem
po prostu podążać tą strategią wygrywającą gracza II. Napotkanie naszego
wcześniejszego ruchu nie jest problemem - wystarczy, zrobić inny losowy ruch.
To oznacza, że gracz I ma strategię wygrywającą - opartą na strategii gracza
II. Sprzeczność, czyli gracz II nie ma strategii wygrywającej, czyli gracz I
ma strategię wygrywającą.

Ten dowód opiera się na kluczowych faktach:
\begin{itemize}
  \item W Hex nie ma remisów
  \item Pola w Hex mogą być tylko w 3 stanach
  \item Ruch dodatkowy w Hex nie może być szkodliwy
\end{itemize}
Najłatwiej ten dowód atakować w innych grach poprzez nasz losowy ruch. Gry takie
jak Rex (odwrócony Hex), charakteryzują się tym, że ruch potrafi być zły dla
gracza.

Kradzież strategii działa też w Chomp. Tam, ponownie zakładamy, że gracz II
ma strategię wygrywającą. Gracz I zjada prawy dolny róg. To jest nieszkodliwe
i potem podąża za strategią gracza II. Zatem mamy sprzeczność - gracz I ma
strategię wygrywającą, a gracz II nie ma strategii wygrywającej.

\section{Metoda ruchów odpowiadających}

Dla niektórych gier, odpowiednią strategią jest ruch odpowiadający. Ogólnie,
można powiedzieć, że ruchem odpowiadającym jest taki ruch, który jest jakimś
odpowiednikiem ruchu gracza przeciwnego. Wtedy strategia staje się funkcją
ruchu przeciwnika.

Na przykład, w kółko i krzyżyk, dla dowolnego rozmiaru planszy i kombinacji,
da się stworzyć parowania pól, które należy wypełnić w momencie kiedy przeciwnik
wypełni sparowane pole. Dla bardziej skomplikowanych wariantów kółka i krzyżyk
da się w ten sposób skonstruować strategię nieprzegrywającą.

\end{document}
